\section{Ethylene và phản ứng của alkene}

\subsection{Cấu tạo của ethylene}

Ethylene có công thức phân tử \ce{C2H4} và công thức cấu tạo:
\[
\ce{CH2=CH2}.
\]
Hai nguyên tử carbon liên kết với nhau bằng một liên kết đôi. Liên kết đôi là nguyên nhân khiến ethylene dễ tham gia phản ứng cộng.

\subsection{Phản ứng cháy}

\[
\ce{C2H4 + 3O2 -> 2CO2 + 2H2O}.
\]

\subsection{Phản ứng với bromine}

Ethylene làm mất màu nước bromine:
\[
\ce{CH2=CH2 + Br2 -> CH2Br-CH2Br}.
\]

\begin{ghinho}{Dấu hiệu nhận biết}
Ở mức THCS, nước bromine là thuốc thử thuận tiện để phân biệt ethylene với methane: ethylene làm mất màu nước bromine, methane không phản ứng trong điều kiện thường.
\end{ghinho}

\subsection{Phản ứng trùng hợp}

Nhiều phân tử ethylene có thể liên kết với nhau tạo polyethylene:
\[
n\ce{CH2=CH2}\longrightarrow \ce{[-CH2-CH2-]_{n}}.
\]
Polyethylene là một polymer được sử dụng rộng rãi trong vật liệu nhựa.

\begin{vidu}
Có hai bình khí không nhãn chứa methane và ethylene. Hãy đề xuất cách nhận biết hai khí bằng một thuốc thử và nêu hiện tượng.
\end{vidu}
\begin{loigiai}
Dẫn từng khí qua nước bromine.

Ethylene phản ứng:
\[
\ce{CH2=CH2 + Br2 -> CH2Br-CH2Br},
\]
làm màu bromine nhạt dần rồi mất màu. Methane không làm mất màu nước bromine trong điều kiện thường.

\textbf{Vậy,} khí làm mất màu nước bromine là ethylene; khí còn lại là methane.
\end{loigiai}

\begin{tuluyen}
Viết phương trình hóa học cho các quá trình sau:
\begin{tasks}(1)
\task ethylene cháy hoàn toàn;
\task ethylene tác dụng với bromine;
\task ethylene trùng hợp tạo polyethylene.
\end{tasks}
\dapso{Xem các phương trình ở phần lí thuyết.}
\end{tuluyen}

\begin{tuluyen}
Cho \SI{0.20}{\mole} ethylene phản ứng hoàn toàn với bromine. Tính số mol bromine cần dùng.
\dapso{\SI{0.20}{\mole}.}
\begin{loigiaiBT}
\[
n_{\ce{C2H4}}=\SI{0.20}{\mole}.
\]

\bangmol{5}
{\ce{C2H4} & + & \ce{Br2} & \rightarrow & \ce{C2H4Br2}}
{0.20 & & 0.20 & & 0.20}

\[
n_{\ce{Br2}}=\SI{0.20}{\mole}.
\]
\textbf{Vậy,} cần \SI{0.20}{\mole} bromine.
\end{loigiaiBT}
\end{tuluyen}
